\par\noindent\textbf{(1)} 稳态方程0=q(E+v×B)−mv/τ。设μ=|q|τ/m、s=sgn(q)，有v\_x−sμB v\_y=sμE\_x、v\_y+sμB v\_x=sμE\_y。求逆得J=nqv=σ0/(1+μ²B²)[[1,sμB],[-sμB,1]]E，其中σ0=nq²τ/m>0；非对角反对称体现Hall响应。
\par\noindent\textbf{(2)} 逆矩阵电阻率ρ=σ\textasciicircum{}\{-1\}=(1/σ0)[[1,−sμB],[sμB,1]]，纵向ρ\_xx=1/σ0与B无关。理想单载流子各向同性Drude模型无本征纵向磁阻，不能因σ\_xx随B变小就直接声称ρ\_xx出现正磁阻，因为实际测量要考虑横向Hall场。
\par\noindent\textbf{(3)} 开路jy=0要求E\_y=ρ\_yx j\_x=sμB j\_x/σ0=B j\_x/(nq)，Hall系数R\_H=E\_y/(j\_xB)=1/(nq)。对q=−e电子，Hall电场与正电荷相反，符号直接反映主导载流子类型；因ρ\_yx随B线性增长，Hall角与μB有关。
\par\noindent\textbf{(4)} 两载流子电子、空穴并存时总σ是两张量相加后再求逆，一般ρ\_xx随B变化而产生正磁阻，也可出现Hall符号翻转。高场量子化、能带各向异性与τ能量依赖超出简单Drude假设。
